% !TeX root = ../MRG41/MRG41_01_poster.tex
%
% 公共加载顺序
% 先载入共享设置与字段命令，再读本期内容，最后应用本期可选排版覆盖。

\usepackage{fontspec, tikz, graphicx, hyperref}

% !TeX root = ../26 09 09 MRG41/MRG41_01_poster.tex
%
% 全项目共用：文字间距、品牌配色、固定标签和校徽路径。
% 根文件声明由构建工具更新，无需随日期前缀手动修改。

% 文字间距：字距单位为 0.01 em；词距为普通空格宽度的倍数。
\newcommand{\GlobalLetterSpacing}{1}   % 1 为略微加宽，负数为收紧
\newcommand{\GlobalWordSpacing}{1.5}   % 1 为原宽度，1.5 为增加 50%
% 词距不改变列距或手动留白；摘要两端对齐时，空格仍会自动伸缩。

% 颜色：只填写六位 HTML 色值，不加 #。
\definecolor{Blue}{HTML}{003C88}        % 论文区深蓝底
\definecolor{Mist}{HTML}{F1F5F8}        % 摘要区、方封面浅蓝底
\definecolor{Ink}{HTML}{102B47}         % 白底上的主信息
\definecolor{Muted}{HTML}{3E566D}       % 白底上的次信息
\definecolor{Pale}{HTML}{D3E4F3}        % 深蓝底上的次信息

% 活动名称与固定标签。
\newcommand{\OrganizerName}{DEEPMACRO LAB}
\newcommand{\EventName}{Macro Reading Group}
\newcommand{\SessionLabel}{SESSION}
\newcommand{\AbstractLabel}{ABSTRACT}

\newcommand{\PresenterLabel}{PRESENTER}
\newcommand{\ScheduleLabel}{SCHEDULE}
\newcommand{\LocationLabel}{LOCATION}
\newcommand{\OnlineLabel}{ONLINE ID}

% 资源路径以每期文件夹为基准；不写本机或 iCloud 的绝对路径。
\newcommand{\LogoPath}{../_config/assets/xmu-emblem.pdf}

% !TeX root = ../MRG41/MRG41_01_poster.tex
%
% 竖版海报：主要字号与距离
% 字号单位为 pt；长度必须带 mm 等单位。行高和总高度由模板自动计算。

% 论文与摘要字号。
\newcommand{\EventNameFontSize}{26}
\newcommand{\PaperTitleFontSize}{36}
\newcommand{\PaperSubtitleFontSize}{20}
\newcommand{\AuthorNameFontSize}{16}
\newcommand{\AuthorAffiliationFontSize}{13}
\newcommand{\AbstractFontSize}{15}
\newcommand{\MetadataFontSize}{14}
\newcommand{\KeywordFontSize}{13}

% 模块 4：同级信息使用同一字号。
\newcommand{\EventPrimarySize}{18}    % 姓名、日期、时间、房间号、会议号
\newcommand{\EventSecondarySize}{15}  % 单位、星期、楼名、会议平台
\newcommand{\EventLabelSize}{12}      % 左侧属性标签

% 整体尺寸与模块留白。
\newlength{\PosterWidth}      \setlength{\PosterWidth}{210mm}
\newlength{\SideMargin}       \setlength{\SideMargin}{16mm}
\newlength{\ModulePadding}    \setlength{\ModulePadding}{10mm}   % 模块上下内边距
\newlength{\ContentGap}       \setlength{\ContentGap}{5mm}       % 内容组间距
\newlength{\EventRowGap}      \setlength{\EventRowGap}{6mm}      % 活动行间附加间距
\newlength{\AuthorRowGap}     \setlength{\AuthorRowGap}{3mm}
\newlength{\AuthorColumnGap}  \setlength{\AuthorColumnGap}{8mm}
\newlength{\LogoWidth}        \setlength{\LogoWidth}{26mm}

% 活动区横向位置；VenueValueX 仅填毫米数，减小会一起左移房间号与会议号。
\newlength{\LabelColumnWidth}  \setlength{\LabelColumnWidth}{44mm}
\newcommand{\VenueValueX}{106}

% !TeX root = ../26 09 09 MRG41/MRG41_02_cover.tex
%
% 公众号封面：比例、字号与安全距离
% 只导出 combined；长图和方图的接缝两侧各加一段底色延展。

% 几何设置：以 100 mm 高母版为基准，其余尺寸随 CoverHeight 同比缩放。
\newcommand{\CoverWideRatio}{2.35}  % 长封面宽高比；方封面固定 1:1
\newcommand{\CoverJoinSafety}{3}    % 每侧安全距离，单位 mm；0 表示关闭
\newcommand{\CoverHeight}{100}      % 整张图的物理高度，单位 mm
\newcommand{\CoverInset}{10}        % 内容距边缘的距离，单位 mm

% 文字层级：主标题与副标题分别控制。
\newcommand{\CoverTitleSize}{36}
\newcommand{\CoverSubtitleSize}{18}
\newcommand{\CoverAuthorSize}{15}
\newcommand{\CoverBrandSize}{11}    % 深蓝区右下角 Lab 与读书会名称
\newcommand{\CoverFooterSize}{10}   % 底部活动信息
\newcommand{\CoverNumberSize}{148}  % 方封面大号期数

% 底部间距与 Lab 标识。
\newcommand{\CoverFooterGap}{2.4em} % 报告人、日期、时间、地点等组间距离
\newcommand{\CoverLogoPath}{../_config/assets/deepmacro-logo.pdf}

% !TeX root = ../MRG41/MRG41_01_poster.tex
%
% 成组填写信息的命令
% 仅简化填写方式；模板继续使用原来的内部宏名。

% 论文发表信息：年份、期刊。
\newcommand{\Publication}[2]{%
    \def\PaperYear{#1}%
    \def\PaperJournal{#2}%
}

% 报告人信息：姓名、单位。
\newcommand{\Presenter}[2]{%
    \def\PresenterName{#1}%
    \def\PresenterAffiliation{#2}%
}

% 活动日期：日期、星期；时间仍独立填写。
\newcommand{\EventDay}[2]{%
    \def\EventDate{#1}%
    \def\EventWeekday{#2}%
}

% 线下地点：楼名、房间号。
\newcommand{\Location}[2]{%
    \def\LocationName{#1}%
    \def\LocationRoom{#2}%
}

% 线上地点：会议平台、会议号。
\newcommand{\OnlineMeeting}[2]{%
    \def\OnlinePlatform{#1}%
    \def\OnlineMeetingID{#2}%
}


% 文件名由本期入口中的 PostNumber 推导，不受文件夹日期前缀影响。
\edef\PostContentFile{MRG\PostNumber_03_content.tex}
\edef\PostLayoutFile{MRG\PostNumber_04_layout.tex}
\input{\PostContentFile}
\InputIfFileExists{\PostLayoutFile}{}{}

% 全文字体：同一套间距应用于 Regular、Medium 与 SemiBold。
\setmainfont{Outfit}[
    Path        = ../_config/assets/fonts/,
    Extension   = .ttf,
    UprightFont = *-Regular,
    BoldFont    = *-SemiBold,
    LetterSpace = \GlobalLetterSpacing,
    WordSpace   = \GlobalWordSpacing
]
\newfontfamily\medium{Outfit-Medium.ttf}[
    Path        = ../_config/assets/fonts/,
    LetterSpace = \GlobalLetterSpacing,
    WordSpace   = \GlobalWordSpacing
]

% 公共排版工具：紧行高字号、同级信息分隔符。
\newcommand{\type}[1]{\fontsize{#1}{#1}\selectfont}
\newcommand{\sep}{\hspace{0.48em}·\hspace{0.48em}}
